% GENERATED FILE. Edit canonical lessons, then run npm run generate:books.
% canonical-source-hash: fnv1a64:eec411f676d3b98a
% canonical-lessons: RU-C289-ogrableniye, RU-C289-grabitel, RU-C289-grabit, RU-C289-napadeniye, RU-C289-napadat

\chapter{Robbery and attack}
\label{ch:ru-robbery-and-attack}
\hlchaptermodality{289}

\begin{chapteropening}
\textbf{By the end of this chapter:} \emph{I can talk about a robbery: a robbery, a robber, robbing, an attack and attacking.}

Complete the last lesson of chapter 289: I can talk about a robbery: a robbery, a robber, robbing, an attack and attacking.
\end{chapteropening}

\section[\texorpdfstring{\ru{ограбление} (ograbléniye)}{ограбление (ograbléniye)}]{\ru{ограбление} (ograbléniye) — a robbery}

\label{lesson:RU-C289-ogrableniye}



\begin{quote}
Before the new one: say the Russian for a pickpocket, then the Russian for a break-in.
\end{quote}

\subsection*{You'll want to know: \ru{ограбление}}
\textbf{\ru{ограбление}} (\emph{ograbléniye}) — \textquotedblleft{}a robbery\textquotedblright{}.

\begin{quote}
\ru{Ограбление} \ru{произошло} \ru{ночью}.
\end{quote}

\begin{quote}
\emph{Ograbléniye proizoshló nóchyu.}
\end{quote}

\begin{quote}
The robbery happened at night.
\end{quote}

Neuter. \emph{\ru{Ограбление} \ru{банка}} — a bank robbery.

\begin{cousinweb}
\emph{\ru{ограбление}} is built on \emph{\ru{ограбить}}, the perfective of \emph{\ru{грабить}} \textquotedblleft{}to rob\textquotedblright{}, which is inherited from Proto-Slavic \emph{grabiti} (EtymDB).
\end{cousinweb}

\begin{grammarlens}[title={What you've built}]
One more word for telling, explaining and describing at B1.
\end{grammarlens}

\subsection*{Your turn}
\begin{itemize}
\raggedright
  \item \emph{Say it:} \emph{ograbléniye}
  \item \emph{Say it:} \emph{ograbléniye}, once more
  \item \emph{Say it:} a sentence of your own with \emph{ograbléniye}
  \item \emph{Recall:} say the Russian for a pickpocket, then the Russian for a break-in, then say \emph{ograbléniye} again
\end{itemize}

\begin{tcolorbox}[breakable,colback=teal!4,colframe=teal!35!black,title={Before you move on}]
What does \ru{ограбление} mean? (\textquotedblleft{}A robbery\textquotedblright{}.) Say it once more.
\end{tcolorbox}
\section[\texorpdfstring{\ru{грабитель} (grabítel)}{грабитель (grabítel)}]{\ru{грабитель} (grabítel) — a robber}

\label{lesson:RU-C289-grabitel}



\begin{quote}
Before the new one: say the Russian for a break-in, then the Russian for a robbery.
\end{quote}

\subsection*{You'll want to know: \ru{грабитель}}
\textbf{\ru{грабитель}} (\emph{grabítel}) — \textquotedblleft{}a robber\textquotedblright{}.

\begin{quote}
\ru{Грабитель} \ru{был} \ru{в} \ru{чёрной} \ru{куртке}.
\end{quote}

\begin{quote}
\emph{Grabítel byl v chyórnoy kúrtke.}
\end{quote}

\begin{quote}
The robber was wearing a black jacket.
\end{quote}

Genitive \emph{\ru{грабителя}}.

\begin{cousinweb}
\emph{\ru{грабитель}} is built on \emph{\ru{грабить}} \textquotedblleft{}to rob\textquotedblright{}, inherited from Proto-Slavic \emph{grabiti} \textquotedblleft{}to grab\textquotedblright{} (EtymDB), with the ending \emph{-\ru{тель}} for a person.
\end{cousinweb}

\begin{grammarlens}[title={What you've built}]
One more word for telling, explaining and describing at B1.
\end{grammarlens}

\subsection*{Your turn}
\begin{itemize}
\raggedright
  \item \emph{Say it:} \emph{grabítel}
  \item \emph{Say it:} \emph{grabítel}, once more
  \item \emph{Say it:} a sentence of your own with \emph{grabítel}
  \item \emph{Recall:} say the Russian for a break-in, then the Russian for a robbery, then say \emph{grabítel} again
\end{itemize}

\begin{tcolorbox}[breakable,colback=teal!4,colframe=teal!35!black,title={Before you move on}]
What does \ru{грабитель} mean? (\textquotedblleft{}A robber\textquotedblright{}.) Say it once more.
\end{tcolorbox}
\section[\texorpdfstring{\ru{грабить} (grábit)}{грабить (grábit)}]{\ru{грабить} (grábit) — to rob}

\label{lesson:RU-C289-grabit}



\begin{quote}
Before the new one: say the Russian for a robbery, then the Russian for a robber.
\end{quote}

\subsection*{You'll want to know: \ru{грабить}}
\textbf{\ru{грабить}} (\emph{grábit}) — \textquotedblleft{}to rob\textquotedblright{}.

\begin{quote}
\ru{Они} \ru{грабили} \ru{магазины} \ru{в} \ru{центре}.
\end{quote}

\begin{quote}
\emph{Oní grábili magazíny v tséntre.}
\end{quote}

\begin{quote}
They robbed shops in the centre.
\end{quote}

Imperfective; perfective \emph{\ru{ограбить}} (\emph{ográbit}): \emph{\ru{его} \ru{ограбили}} — he was robbed. Present \emph{\ru{граблю}}, \emph{\ru{грабишь}}.

\begin{cousinweb}
\emph{\ru{грабить}} is inherited from Proto-Slavic \emph{grabiti} (EtymDB), usually glossed in its old sense \textquotedblleft{}to grab, to rake together\textquotedblright{}.
\end{cousinweb}

\begin{grammarlens}[title={What you've built}]
One more word for telling, explaining and describing at B1.
\end{grammarlens}

\subsection*{Your turn}
\begin{itemize}
\raggedright
  \item \emph{Say it:} \emph{grábit}
  \item \emph{Say it:} \emph{grábit}, once more
  \item \emph{Say it:} a sentence of your own with \emph{grábit}
  \item \emph{Recall:} say the Russian for a robbery, then the Russian for a robber, then say \emph{grábit} again
\end{itemize}

\begin{tcolorbox}[breakable,colback=teal!4,colframe=teal!35!black,title={Before you move on}]
What does \ru{грабить} mean? (\textquotedblleft{}To rob\textquotedblright{}.) Say it once more.
\end{tcolorbox}
\section[\texorpdfstring{\ru{нападение} (napadéniye)}{нападение (napadéniye)}]{\ru{нападение} (napadéniye) — an attack}

\label{lesson:RU-C289-napadeniye}



\begin{quote}
Before the new one: say the Russian for a robber, then the Russian for to rob.
\end{quote}

\subsection*{You'll want to know: \ru{нападение}}
\textbf{\ru{нападение}} (\emph{napadéniye}) — \textquotedblleft{}an attack\textquotedblright{}.

\begin{quote}
\ru{Нападение} \ru{было} \ru{неожиданным}.
\end{quote}

\begin{quote}
\emph{Napadéniye býlo neozhídannym.}
\end{quote}

\begin{quote}
The attack was unexpected.
\end{quote}

\emph{\ru{Нападение} \ru{на}} + accusative — an attack on.

\begin{cousinweb}
\emph{\ru{нападение}} is built on \emph{\ru{нападать}} \textquotedblleft{}to attack\textquotedblright{}, from \emph{\ru{на}-} \textquotedblleft{}on\textquotedblright{} and \emph{\ru{падать}} \textquotedblleft{}to fall\textquotedblright{}, inherited from Proto-Slavic \emph{padati} (EtymDB): a falling upon.
\end{cousinweb}

\begin{grammarlens}[title={What you've built}]
One more word for telling, explaining and describing at B1.
\end{grammarlens}

\subsection*{Your turn}
\begin{itemize}
\raggedright
  \item \emph{Say it:} \emph{napadéniye}
  \item \emph{Say it:} \emph{napadéniye}, once more
  \item \emph{Say it:} a sentence of your own with \emph{napadéniye}
  \item \emph{Recall:} say the Russian for a robber, then the Russian for to rob, then say \emph{napadéniye} again
\end{itemize}

\begin{tcolorbox}[breakable,colback=teal!4,colframe=teal!35!black,title={Before you move on}]
What does \ru{нападение} mean? (\textquotedblleft{}An attack\textquotedblright{}.) Say it once more.
\end{tcolorbox}
\section[\texorpdfstring{\ru{нападать} (napadát)}{нападать (napadát)}]{\ru{нападать} (napadát) — to attack}

\label{lesson:RU-C289-napadat}



\begin{quote}
Before the new one: say the Russian for to rob, then the Russian for an attack.
\end{quote}

\subsection*{You'll want to know: \ru{нападать}}
\textbf{\ru{нападать}} (\emph{napadát}) — \textquotedblleft{}to attack\textquotedblright{}.

\begin{quote}
\ru{Собака} \ru{напала} \ru{на} \ru{почтальона}.
\end{quote}

\begin{quote}
\emph{Sobáka napála na pochtalóna.}
\end{quote}

\begin{quote}
The dog attacked the postman.
\end{quote}

Imperfective; perfective \emph{\ru{напасть}} (\emph{napást}), past \emph{\ru{напал}}. It takes \emph{\ru{на}} + accusative.

\begin{cousinweb}
\emph{\ru{нападать}} is built on \emph{\ru{на}-} \textquotedblleft{}on\textquotedblright{} and \emph{\ru{падать}} \textquotedblleft{}to fall\textquotedblright{}, which is inherited from Proto-Slavic \emph{padati} (EtymDB): to fall upon.
\end{cousinweb}

\begin{grammarlens}[title={What you've built}]
One more word for telling, explaining and describing at B1.
\end{grammarlens}

\subsection*{Your turn}
\begin{itemize}
\raggedright
  \item \emph{Say it:} \emph{napadát}
  \item \emph{Say it:} \emph{napadát}, once more
  \item \emph{Say it:} a sentence of your own with \emph{napadát}
  \item \emph{Recall:} say the Russian for to rob, then the Russian for an attack, then say \emph{napadát} again
\end{itemize}

\begin{tcolorbox}[breakable,colback=teal!4,colframe=teal!35!black,title={Before you move on}]
What does \ru{нападать} mean? (\textquotedblleft{}To attack\textquotedblright{}.) Say it once more.
\end{tcolorbox}
